\section{Complex Scaling and the ABC Theorem}
\subsection{Spectral Theorem}
A Hilbert space is a complete inner product space:
\begin{itemize}
    \item Inner product space: a vector space equipped with an inner product (notion of distance, length, and angles and a way to add states together).
    \item Completeness: every Cauchy sequence $(a_i)_{i\in I}$ in $\mathcal{H}$ has a limit in $\mathcal{H}$ (notion of continuity).
\end{itemize}
The adjoint of an operator $H$ in an inner product space $S$ is the operator $H^\dagger$ s.t. $\langle H^\dagger \psi, \phi\rangle = \langle \psi, H\phi\rangle$ $\forall$ $\psi, \phi$ $\in S$. If an operator is self-adjoint, then $H^\dagger=H$ or $\langle H \psi, \phi\rangle = \langle \psi, H\phi\rangle$. Generally, we work with the Hilbert space $\mathcal{H}\subseteq L^2(\mathbb{R}^2)$ equipped with the Hermitian inner product. In this case, the Hamiltonian is self-adjoint (in general, we assume that the Hamiltonian in any physical Hilbert space is self-adjoint). Self-adjoint operators under the Hermitian inner product have real eigenvalues. The spectral theorem guarantees that, for any self-adjoint operator $H$, there exists a projection-valued measure $P(\lambda)$ (projection-valued measure guarantees that $dP(\lambda)$ is an orthogonal projector with respect to the inner product) such that
\[
    H = \int_\mathbb{R} \lambda P(\lambda) \text{ and } I = \int_\mathbb{R} dP(\lambda).
\]
For a general Hamiltonian, the spectrum is usually of the form $\sigma(H) = \{E_1, E_2, \dots, E_n\}\;\cup\; [E_\text{th},\infty)$ (the union of the discrete spectrum and continuous spectrum), and $P_\text{cont}|n\rangle = 0$ (the projection of a bound state onto the continuum subspace is null and visa vera). Then,
\[
    H = \sum_n E_n |n\rangle\langle n| + \int_{E_\text{th}}^\infty \; dE\;E\; |E\rangle\langle E| \text{ and } I = \sum_n |n\rangle\langle n| + \int_{E_\text{th}}^\infty dE |E\rangle \langle E|
\]
We rely heavily on the self-adjoint-ness of $H$ (and all other operators representing physical observables) and the spectral theorem in quantum mechanics. Once we complex scale the Hamiltonian, it is no longer self-adjoint with respect to the Hermitian inner product, and many of our assumptions break down. However, if we equip our Hilbert space with the c-product instead of the Hermitian inner product, the self-adjoint-ness of $H$ and the spectral theorem are recovered: $H=H^\dagger$ or $(H\psi, \phi) = (\psi, H\phi)$ $\forall$ $\psi, \phi\in \mathcal{H}$.

\subsection{General Idea Behind CS}

% brief introduction into complex scaling and how resonances are affected
In complex scaling, one can show that making the substituion $r\rightarrow re^{i\theta}$ can make resonances appear bound. This is analytically continuing the Gamow state to the complex plane and evaluating it on a complex $r$-ray. Because of the $e^{i\theta}$, this can damp the diverging oscillations (for suitable choice of $\theta$) and make the Gamow state $L^2$ normalizable.

The ABC theorem says that: if the potential is dilation-analytic (the analytic continuation of the potential is analytic in a specific $\theta$ wedge), then the complex scaled Hamiltonian is analytic in that wedge, the continuum is rotated to the position of the ray, and resonances become bound states (in that specific $\theta$ wedge) with the same resonance energy as before. This has huge computational advantages because it allows us to find resonances through bound state methods, and is usually easy to implement computationally.

Because the complex scaled Hamiltonian is no longer normalizable, the Hermitian inner product is no longer a feasible definition for the inner product (one can normalize according to any convention one wants, and you can certainly recover the same physics if you continue to use the Hermitian inner product, but you have to be extremely careful to recognize assumptions you'll probably naturally make even though they no longer hold). Mainly, eigenvectors of the Hamiltonian are no longer naturally orthonormal wrt the Hermitian inner product, which makes insertions of identity unfortunate to work with. Defining an inner product as the c-product fixes these issues: left and right eigenvectors become identical again, the dual space structure naturally implies the left-eigenvector structure, and eigenvectors of the Hamiltonian are once again orthonormal).

\subsection{\texorpdfstring{Proof that most observables are $\phi$ independent}{Proof that most observables are phi independent}}\label{sssection:obs_phi_indep}
One can further show that, for any observable $A$ that is dilation-analytic and doesn't blow up pathologically (falls off faster than a specific exponential rise), expectation values like $\langle A_\theta \rangle$ are $\theta$ independent. I.e., the expectation values are regularized observables. I'll put this proof here shortly. The general idea is to show that $\langle A_{\theta+\delta\theta}\rangle = \langle A_{\theta}\rangle$ (where $\theta$ and $\theta+\delta\theta$ are in the good sector) by showing that the integrals end up being equivalent because $A$ is dilation analytic in the wedge, so there is no residue accumulated by integrating over the $\delta \theta$ ray.
